% =====================================================================
%  PLANTILLA BEAMER SDS - Cátedra 72.25 Simulación de Sistemas
%  Estructura validada según GuiaPresentaciones.md y feedback docente.
%
%  Compilación:
%   - Para entrega en PDF (muestra fotograma estático + link YouTube):
%       pdflatex presentacion.tex
%   - Para generar .pptx en vivo con GIFs embebidos:
%       pdflatex -jobname=presentacion_vivo "\def\modo{vivo}\input{presentacion}"
%       py scripts/build_pptx.py
% =====================================================================
\documentclass[aspectratio=169]{beamer}

\usetheme{Boadilla}
\useoutertheme{miniframes}
\usecolortheme{whale}

\usepackage[utf8]{inputenc}
\usepackage[T1]{fontenc}
\usepackage[spanish,es-noquoting]{babel}
\usepackage{amsmath,amssymb,graphicx,booktabs}
\usepackage{tikz}
\usetikzlibrary{arrows.meta,positioning,fit,backgrounds,calc}

% Guía 3.5: Vectores en negrita sin itálica, escalares en itálica sin negrita.
\newcommand{\vect}[1]{\mathbf{#1}}

\setbeamertemplate{navigation symbols}{}
\setbeamertemplate{footline}[frame number]
\setbeamercovered{transparent}

% Diapositiva divisoria automática por sección (Guía 1.13)
\AtBeginSection[]{%
  \begin{frame}[plain]
    \vfill
    \begin{beamercolorbox}[sep=8pt,center]{title}
      \usebeamerfont{title}\insertsection
    \end{beamercolorbox}
    \vfill
  \end{frame}
}

% Bloque de parámetros al costado de una figura (Guía 1.7)
\newcommand{\params}[1]{{\footnotesize #1}}

% --- Modo de compilación: entrega (default) vs. vivo ------------------
\providecommand{\modo}{entrega}
\def\modovivo{vivo}

% Color marcador del hueco de animación para build_pptx.py (RGB 255, 0, 255)
\definecolor{huecoanim}{RGB}{255,0,255}
\newcommand{\soloentrega}[1]{\ifx\modo\modovivo\else#1\fi}

% Macro para insertar animaciones
% #1: Fotograma PNG, #2: Título/subtítulo, #3: Link a video
\newcommand{\animacion}[3]{%
  \centering
  \ifx\modo\modovivo
    \tikz\fill[huecoanim] (0,0) rectangle (0.48\textheight,0.48\textheight);%
  \else
    \includegraphics[height=0.48\textheight]{#1}%
  \fi
  \\[2pt]
  {\scriptsize #2}\\
  \soloentrega{{\tiny\color{blue}\url{#3}}}%
}

% --- Metadatos de la Presentación ------------------------------------
\title[Título Corto]{Título Completo del Trabajo Práctico}
\subtitle{Subtítulo o Modelo Físico Estudiado}
\author[Grupo G]{Integrante 1 \and Integrante 2 \and Integrante 3}
\institute[ITBA]{Grupo G --- Comisión S \\ Instituto Tecnológico de Buenos Aires}
\date{72.25 Simulación de Sistemas --- Trabajo Práctico N\textdegree{}X --- 2026}

\begin{document}

% Portada
\begin{frame}[plain]
  \titlepage
\end{frame}

% =====================================================================
\section{Introducción}
% =====================================================================

\begin{frame}{Sistema Real}
  \begin{columns}[c]
    \begin{column}{0.55\textwidth}
      \centering
      % \includegraphics[height=0.65\textheight,keepaspectratio]{sistema_real.png}
    \end{column}
    \begin{column}{0.43\textwidth}
      \begin{itemize}
        \item Fenómeno físico observado en la naturaleza o ingeniería.
        \item Características dinámicas relevantes.
        \item Motivación del estudio por simulación computacional.
      \end{itemize}
    \end{column}
  \end{columns}
\end{frame}

\begin{frame}{Modelo Matemático}
  \begin{itemize}
    \item Leyes físicas que gobiernan el sistema:
      \begin{equation}
        \vect{v}_i(t + \Delta t) = \dots
      \end{equation}
    \item Ecuaciones de conservación e interacciones generales.
    \item Sin introducir discretizaciones numéricas particulares todavía.
  \end{itemize}
\end{frame}

% =====================================================================
\section{Implementación}
% =====================================================================

\begin{frame}{Arquitectura del Motor de Simulación}
  \begin{columns}[c]
    \begin{column}{0.60\textwidth}
      % Diagrama de arquitectura / diagrama UML / ciclo principal
    \end{column}
    \begin{column}{0.38\textwidth}
      \begin{itemize}
        \item Algoritmo del motor (C++).
        \item Estructura de datos (e.g. Cola de eventos, grilla espacial).
        \item \textbf{Nota cátedra:} Solo el motor; post-proceso excluido.
      \end{itemize}
    \end{column}
  \end{columns}
\end{frame}

% =====================================================================
\section{Simulaciones}
% =====================================================================

\begin{frame}{Sistema Simulado}
  \begin{columns}[c]
    \begin{column}{0.55\textwidth}
      % Esquema geométrico con parámetros acotados
    \end{column}
    \begin{column}{0.43\textwidth}
      \params{
        Geometría: $L \times W$\\[4pt]
        Partículas: $N = 100$, $r = 1{,}75\times10^{-2}$ m\\[4pt]
        Condiciones de contorno: periódicas / rígidas\\[4pt]
        Realizaciones: 100 por punto
      }
    \end{column}
  \end{columns}
\end{frame}

\begin{frame}{Definición de Observables}
  \begin{itemize}
    \item Definición matemática estricta de las métricas evaluadas:
      \begin{equation}
        \langle \mathcal{O} \rangle = \frac{1}{M} \sum_{k=1}^M \mathcal{O}_k
      \end{equation}
    \item Barra de error: error estándar de la media ($\sigma / \sqrt{M}$), declarada aquí para toda la presentación.
  \end{itemize}
\end{frame}

% =====================================================================
\section{Resultados}
% =====================================================================

% Tríada: 1. Dinámica (animación) -> 2. Evolución temporal -> 3. Input vs Observable
\begin{frame}{Parámetro $X$: Dinámica del Sistema}
  \vfill
  \begin{columns}[c]
    \begin{column}{0.45\textwidth}
      \animacion{frame_caso_a.png}{Caso Extremo A (bajo $X$)}{https://youtu.be/PENDIENTE}
    \end{column}
    \begin{column}{0.45\textwidth}
      \animacion{frame_caso_b.png}{Caso Extremo B (alto $X$)}{https://youtu.be/PENDIENTE}
    \end{column}
    \begin{column}{0.10\textwidth}
      \params{$t = 10$ s}
    \end{column}
  \end{columns}
  \vfill
\end{frame}

\begin{frame}{Parámetro $X$: Evolución Temporal}
  \begin{columns}[c]
    \begin{column}{0.72\textwidth}
      % \centering\includegraphics[width=\linewidth,height=0.80\textheight,keepaspectratio]{evolucion_temporal.png}
    \end{column}
    \begin{column}{0.26\textwidth}
      \params{
        Verificación de estado estacionario o tiempo característico.
      }
    \end{column}
  \end{columns}
\end{frame}

\begin{frame}{Parámetro $X$ vs Observable}
  \begin{columns}[c]
    \begin{column}{0.70\textwidth}
      % \centering\includegraphics[width=\linewidth,height=0.80\textheight,keepaspectratio]{input_vs_observable.png}
    \end{column}
    \begin{column}{0.28\textwidth}
      \params{
        Símbolos: promedio sobre 100 semillas.\\[4pt]
        Barras: error estándar.\\[4pt]
        Línea: guía visual (sin interpolación espuria).
      }
    \end{column}
  \end{columns}
\end{frame}

% =====================================================================
\section{Conclusiones}
% =====================================================================

\begin{frame}{Conclusiones}
  \vfill
  \begin{itemize}
    \item Conclusión 1: Respaldada directamente por los datos presentados.
    \vspace{0.6em}
    \item Conclusión 2: Comparación cuantitativa entre configuraciones ($A$ vs $B$).
    \vspace{0.6em}
    \item Conclusión 3: Tendencia física o leyes de escala verificadas.
  \end{itemize}
  \vfill
\end{frame}

% Cierre
\begin{frame}[plain]
  \vfill
  \centering
  \Huge \textbf{Muchas Gracias}
  \vfill
\end{frame}

\end{document}
